% ===========================================================================
% METHOD --- execution-protocol specification.
% This section defines how the trusted controller and untrusted accelerator
% execute inference and LoRA adaptation.
% ===========================================================================

% The local TeX installation used by this artifact does not provide the
% algorithm/algorithmic packages. This small, package-free environment retains
% AAAI's two-column layout and gives the pseudocode an independent counter.
\newcounter{methodalgorithm}
\newenvironment{methodalgorithm}[1]{%
  \par\medskip
  \refstepcounter{methodalgorithm}%
  \small\hrule\vspace{3pt}
  \noindent\textbf{Algorithm \themethodalgorithm: #1}\par\vspace{3pt}
}{%
  \vspace{3pt}\hrule
  \par\medskip
}

\section{Method}
\label{sec:method}

\subsection{System Overview}
\label{subsec:method-overview}

The execution has two administrative phases and two operating modes. During
offline preparation, the model provider installs the private parameters
$\params$ in the trusted controller $\ctrl$. The controller samples the masks
needed by the configured call schedule, validates every projection boundary,
and records the execution order. In online operation, $\ctrl$ constructs the
transformed operands for each call and alternates between trusted operators over
plaintext transients and calls to the untrusted accelerator $\acc$ over tensors
in $\Sobs$. In inference mode, one prompt prefill is followed by one-token
decode calls until generation terminates. In adaptation mode, the same masked
linear interface is extended with separate, rank-padded LoRA factors; gradient
recovery and the optimizer interface remain in $\ctrl$.

The trusted controller owns the token embedding lookup, all masks and inverse
masks, pads, recovery maps, normalization gains, plaintext residual state,
plaintext attention projections after recovery, recovered logits, token
selection, the plaintext LoRA factors, recovered adapter gradients, and
optimizer state. The accelerator owns no plaintext model or runtime state. It
stores the transformed base projections, transformed LoRA factors when active,
and the session's masked key/value cache, and it evaluates the requested dense
linear algebra and the compatible attention and activation kernels. A boundary
call sends only the transformed operands and public shapes needed for that
call; its response is a transformed tensor that $\ctrl$ either recovers or
passes to a subsequent compatible accelerator operation.

\begin{figure*}[t]
\centering
\fbox{\begin{minipage}{0.94\textwidth}
\centering
\textbf{Offline preparation}\quad
$\params,(\Amat,\Bmat),\maskset\in\ctrl$
$\longrightarrow$
$\{\text{validated shapes, call order, fresh schedule}\}$ in $\ctrl$

\vspace{5pt}
$\usr:x$
$\longrightarrow$
$\boxed{\ctrl:\ \text{embed, normalize, mask/pad}}$
$\longrightarrow$
$\boxed{\acc:\ \text{masked linear/attention/activation}}$
$\longrightarrow$
$\boxed{\ctrl:\ \text{recover, residual update, logits, token}}$

\vspace{5pt}
Prefill initializes $\{\Kt,\Vt\}_{\ell=1}^{L}$; each decode call appends one
masked key/value row per layer and returns one recovered next-token decision.
LoRA adaptation replaces a delegated linear call by a base branch plus the
separate $\At\Bt$ branch, followed by masked backward and a trusted update.
\end{minipage}}
\caption{Overview placeholder for the execution lifecycle. The final artwork
will replace this boxed protocol sketch without changing the depicted boundary
or message order.}
\label{fig:method-overview}
\end{figure*}

\begin{methodalgorithm}{High-level controller--accelerator execution}
\label{alg:method-overview}
\textbf{Inputs:} private token sequence $x$; private parameters $\params$;
optional private adapter $(\Amat,\Bmat)$; public architecture $\arch$.\par
\textbf{Outputs:} generated token sequence, or updated private adapter.\par
\textbf{Trusted state:} $\params,\maskset,\Tpad,\Recover$; optional adapter and
optimizer state; plaintext residuals, logits, and gradients.\par
\textbf{Accelerator state:} transformed operators and optional transformed
adapter factors; masked KV cache; current transformed activations.\par
\textbf{Communication:} calls labelled \textsc{Send} and \textsc{Return}.\par
\begin{enumerate}
\setlength{\itemsep}{0pt}\setlength{\parskip}{0pt}\setlength{\parsep}{0pt}
\item $\ctrl$ installs the private model and optional adapter, validates the
call plan, and initializes its fresh-mask schedule.
\item $\ctrl$ embeds the prompt, performs the trusted boundary operations, and
\textsc{Send}s the first transformed activation.
\item $\acc$ executes the masked prefill and \textsc{Return}s transformed
layer outputs and masked cache entries as requested.
\item $\ctrl$ recovers the final logits and selects the next token.
\item While generation continues, repeat the one-token boundary sequence,
reusing the session KV masks and appending to the masked cache.
\item If adapting, execute masked forward and backward calls, recover the
adapter gradients in $\ctrl$, update the private factors, and refresh their
transformed images.
\item Erase pass-local plaintexts and return only the selected output.
\end{enumerate}
\end{methodalgorithm}

Figure~\ref{fig:method-overview} shows the lifecycle, and
Algorithm~\ref{alg:method-overview} fixes its message order. The remaining
subsections specify each step.

\subsection{Runtime Protocol}
\label{subsec:runtime-protocol}

\paragraph{Offline phase.}
For every linear map with private weight
$\W\in\mathbb{R}^{\dinX\times\doutX}$, $\ctrl$ records its input and output
dimensions and the mask family required by the surrounding operators. It
validates the attention head grouping, RoPE layout, cache layout, normalization
locations, and residual order. For each adapted projection it creates the
rank-padded controller representation and initializes its private optimizer
state. It may precompute a trusted schedule of distinct call-local mask and pad
material; this schedule remains inside $\ctrl$. No plaintext or transformed
model operand is installed on $\acc$ in this phase.

\paragraph{Runtime phase.}
The user supplies token identifiers to $\ctrl$. The controller performs the
embedding lookup and holds the resulting plaintext hidden state $\X$. At every
delegated linear boundary, $\ctrl$ samples a fresh $\Nin$, $\Nout$, and
$\Tpad$ for the call, or consumes the corresponding precomputed trusted
schedule entry. It constructs $\Wt=\Nin^{-1}\W\Nout$, the transformed bias,
and $\Ct=\Tpad\W\Nout$, then sends these operands with
$\Xt=(\X-\Tpad)\Nin$. The accelerator evaluates the requested multiplication,
adds the transformed bias and compensation, and returns $\Yt$. The controller
applies $\Recover$ when the next operator is trusted; when the next operation
accepts the same transformed basis, $\Yt$ remains on the accelerator until that
operation finishes. Call-local transformed operands are released after their
consumer finishes.

\paragraph{Communication and boundary crossings.}
A standard linear call has one controller-to-accelerator message containing
$\Xt,\Wt,\bias\Nout,\Ct$ and public shape metadata, followed by one
accelerator-to-controller message containing $\Yt$. Attention uses the same
messages for the $Q$, $K$, $V$, and output projections; the accelerator
additionally retains $\Kt$ and $\Vt$ between calls. A nonlinear boundary returns the current masked
activation to $\ctrl$, which recovers it, applies the trusted operator, and
starts the next linear call with a new mask and pad. Final masked logits cross
from $\acc$ to $\ctrl$ once per prefill or decode call; selected tokens do not
form an accelerator request.

\paragraph{Freshness schedule.}
Boundary masks, rank mixers, and additive pads are fresh for each delegated
call. A decode schedule may precompute these call-local secrets inside
$\ctrl$; each schedule entry is consumed once. The key and value masks
$\NK,\NV$ are different: they are sampled once per generation session and per
layer, then held fixed through prefill and every decode append. The query mask
is derived from the corresponding key mask for each query head. Termination
destroys the cache masks and all unused call-local schedule entries. A new
generation session samples a new set.

\paragraph{Tensor ownership.}
The private parameters, masks and inverses, raw pads, plaintext residuals,
recovered $Q/K/V$, recovered logits, adapter factors, real-rank gradients, and
optimizer state are controller-owned. The transformed base and adapter
operators, $\Xt$, intermediate masked activation branches, $\Yt$, and masked
KV cache are accelerator-held. Architecture, tensor shapes, sequence and cache
lengths, and $\rpad$ are public metadata. Compensation tensors are produced by
$\ctrl$ and consumed by $\acc$; they are not recovery state. The recovery maps
and inverse masks are never installed on $\acc$.

\subsection{Linear Execution}
\label{subsec:linear-execution}

Consider one single-sequence linear call in row-vector convention. Its
controller input is $\X\in\mathbb{R}^{n\times\dinX}$, its private weight is
$\W\in\mathbb{R}^{\dinX\times\doutX}$, its optional public bias is
$\bias\in\mathbb{R}^{\doutX}$, and its plaintext output slot is
$\Y\in\mathbb{R}^{n\times\doutX}$. Batched execution applies the same maps to
each batch item. The controller samples
$\Nin,\Nin^{-1}\in\mathbb{R}^{\dinX\times\dinX}$,
$\Nout,\Nout^{-1}\in\mathbb{R}^{\doutX\times\doutX}$, and
$\Tpad\in\mathbb{R}^{n\times\dinX}$. Each mask is invertible; operator-specific
boundaries may restrict it to an orthogonal or permutation family.

The controller constructs the right-action input and transformed operator
\begin{align}
\Xt&=(\X-\Tpad)\Nin
      \in\mathbb{R}^{n\times\dinX},\\
\Wt&=\Nin^{-1}\W\Nout
      \in\mathbb{R}^{\dinX\times\doutX},
\label{eq:method-linear-transform}
\end{align}
together with
$\bias\Nout\in\mathbb{R}^{\doutX}$ and the compensation
\begin{equation}
\Ct=\Tpad\W\Nout\in\mathbb{R}^{n\times\doutX}.
\label{eq:method-linear-compensation}
\end{equation}
For a broadcast pad, the implementation sends
$\Tpad\Nin\in\mathbb{R}^{\dinX}$ and
$\Ct\in\mathbb{R}^{\doutX}$ and broadcasts both over the $n$ rows. The pad is
local to this linear boundary and is compensated before its output is consumed
by RoPE, softmax, normalization, or an activation.

The accelerator receives or references only the four transformed operands and
computes
\begin{equation}
\Yt=\Xt\Wt+\bias\Nout+\Ct
\in\mathbb{R}^{n\times\doutX}.
\label{eq:method-linear-accelerator}
\end{equation}
If the result returns to a trusted operator, $\ctrl$ applies
\begin{equation}
\Recover(\Yt)=\Yt\Nout^{-1}
\in\mathbb{R}^{n\times\doutX}.
\label{eq:method-linear-recovery}
\end{equation}
If the next accelerator operator is defined in the $\Nout$ basis, $\Yt$ stays
on $\acc$ and that operator's input fold uses $\Nout^{-1}$ directly. Thus
$\ctrl$ owns $\X,\W,\Tpad,\Nin^{\pm1},\Nout^{\pm1},\Recover$, while $\acc$
owns the runtime copies of $\Xt,\Wt,\bias\Nout,\Ct,\Yt$.

\subsection{Attention Execution}
\label{subsec:attention-execution}

At layer $\ell$, the normalized input has shape
$\X\in\mathbb{R}^{n\times d}$. The four private projection matrices are
$\W_Q\in\mathbb{R}^{d\times h d_h}$,
$\W_K\in\mathbb{R}^{d\times h_{kv}d_h}$,
$\W_V\in\mathbb{R}^{d\times h_{kv}d_h}$, and
$\W_O\in\mathbb{R}^{h d_h\times d}$. Each projection first uses the linear
protocol of Section~\ref{subsec:linear-execution}. Its recovered outputs are
reshaped into per-head matrices
$\Qmat\in\mathbb{R}^{n\times d_h}$ for each of the $h$ query heads and
$\K,\Vmat\in\mathbb{R}^{n\times d_h}$ for each of the $h_{kv}$ key/value
heads.

The wired path applies $\RoPE$ to the recovered $\Qmat$ and $\K$ inside
$\ctrl$, using the public absolute token positions. The controller then
right-multiplies each rotated query by
$\NQ\in\mathbb{R}^{d_h\times d_h}$ and each rotated key by
$\NK\in\mathbb{R}^{d_h\times d_h}$, with
$\NQ\NK^{\mathsf T}=\mathbf{I}_{d_h}$. Values are right-multiplied by
$\NV\in\mathbb{R}^{d_h\times d_h}$. The resulting
$\Qt,\Kt,\Vt\in\mathbb{R}^{n\times d_h}$ are accelerator tensors.

For grouped-query attention, each key/value head is shared by exactly
$h/h_{kv}$ query heads. The controller installs one $\NK$ and one $\NV$ for
that shared key/value head and derives the query mask of every query head in
the group from the same $\NK$. Multi-query attention is the case
$h_{kv}=1$; all query heads use the one session key mask through their derived
query masks. No mask mixes distinct attention heads.

During prefill, $\acc$ appends the rotated masked keys and masked values to an
initially empty per-layer cache. The cache tensors have shape
$\Kt,\Vt\in\mathbb{R}^{n\times d_h}$ per key/value head. The accelerator
repeats each cached key/value head across its query-head group, forms the
$n\times n$ score matrix for each query head, applies the causal bias and
softmax, and multiplies the resulting $n\times n$ probability matrix by the
repeated $\Vt\in\mathbb{R}^{n\times d_h}$. The $h$ outputs are concatenated
into an $n\times h d_h$ matrix, and the output projection maps it to an
$n\times d$ transformed residual update. The controller recovers that update
before the residual addition in the wired path.

During decode position $t$, the query, new key, and new value each have one
token row. The accelerator appends the new
$\Kt,\Vt\in\mathbb{R}^{1\times d_h}$ to the existing
$t\times d_h$ per-head cache, yielding $(t+1)\times d_h$ cache matrices. It
forms a $1\times(t+1)$ score row, applies softmax over the cache-length axis,
and produces a $1\times d_h$ value aggregate. The same session $\NK$ and
$\NV$ are used for every append; call-local linear masks and pads are freshly
sampled for the projections. The output projection, residual update, following
normalization, and LM head then execute in the same order as prefill.

\subsection{Nonlinear Execution}
\label{subsec:nonlinear-execution}

The wired trusted-boundary path executes embedding lookup, RMSNorm (including
the private gain $\rmsgain$), final normalization, RoPE, recovery, residual
addition, logit recovery, and token selection inside $\ctrl$. Before either
normalization in a decoder block, $\ctrl$ recovers the current hidden state,
applies the normalization to obtain a plaintext transient, and starts the next
$Q/K/V$ or MLP projection boundary with a fresh input mask and additive pad.
After the attention output projection returns, $\ctrl$ recovers its
$n\times d$ output and adds it to the controller-held residual. It repeats the
same sequence for the MLP output.

The attention softmax and value aggregation execute on $\acc$ over the masked
$Q/K/V$ representation described above. The SwiGLU activation also executes on
$\acc$: $\ctrl$ chooses one channel permutation $\Perm$ for both the gate and
up projections, and their accelerator outputs are
$n\times m$ matrices in that shared permuted basis. The accelerator applies
$\SiLU$ to the gate branch, multiplies it elementwise with the up branch, and
uses a down projection whose input fold contains $\Perm^{-1}$. GELU and ReLU
operators, when selected by the model architecture, use the same permutation
boundary. Additive pads are compensated at the preceding linear outputs and
are not carried into any activation or normalization core.

Plaintext activations produced inside $\ctrl$ are pass-local. A trusted
nonlinear call consumes the recovered input, produces the next plaintext
transient, and releases only the newly masked input of the following linear
boundary. Accelerator-side activations remain in their operator-compatible
permutation or attention basis until their folded output projection returns a
transformed residual update.

\subsection{LoRA Adaptation}
\label{subsec:lora-adaptation}

For an adapted projection, $\ctrl$ owns
$\Amat\in\mathbb{R}^{\dinX\times\rtrue}$ and
$\Bmat\in\mathbb{R}^{\rtrue\times\doutX}$. It first constructs
$\Apad\in\mathbb{R}^{\dinX\times\rpad}$ and
$\Bpad\in\mathbb{R}^{\rpad\times\doutX}$, with the real factors occupying
the first $\rtrue$ columns and rows. The remaining blocks are either a random
down block paired with a zero up block, or pairs
$[\Adummy,\Adummy]$ and $[\Bdummy;-\Bdummy]$. The scaling used by all forward
and backward calls is $\alpha/\rtrue$; padding does not change this
denominator.

For one call, $\ctrl$ samples
$\Umix,\Umix^{-1}\in\mathbb{R}^{\rpad\times\rpad}$ and constructs
\begin{align}
\At&=\Nin^{-1}\Apad\Umix
     \in\mathbb{R}^{\dinX\times\rpad},\\
\Bt&=\Umix^{-1}\Bpad\Nout
     \in\mathbb{R}^{\rpad\times\doutX}.
\label{eq:method-lora-factors}
\end{align}
The accelerator evaluates the base and adapter branches without merging their
factor representation:
\begin{equation}
\Yt=\Xt\Wt+
\frac{\alpha}{\rtrue}(\Xt\At)\Bt+
\bias\Nout+\Ct,
\label{eq:method-lora-forward}
\end{equation}
where $\Xt\At\in\mathbb{R}^{n\times\rpad}$,
$(\Xt\At)\Bt\in\mathbb{R}^{n\times\doutX}$, and $\Ct$ contains the base and
adapter pad compensations. The factors remain separate accelerator operands
throughout the adaptation call.

For backward execution, $\ctrl$ holds the upstream gradient
$\Gmat\in\mathbb{R}^{n\times\doutX}$ and sends
$\Gt=\Gmat\Nout^{-\mathsf T}\in\mathbb{R}^{n\times\doutX}$. The accelerator
computes
\begin{align}
d\At &=\frac{\alpha}{\rtrue}\Xt^{\mathsf T}
       (\Gt\Bt^{\mathsf T})
       \in\mathbb{R}^{\dinX\times\rpad},\\
d\Bt &=\frac{\alpha}{\rtrue}(\Xt\At)^{\mathsf T}\Gt
       \in\mathbb{R}^{\rpad\times\doutX}.
\label{eq:method-lora-backward}
\end{align}
It returns these two masked gradients. The controller maps them to the padded
factor coordinates as
\begin{align}
d\Apad &=\Nin^{-\mathsf T}d\At\Umix^{\mathsf T}
          \in\mathbb{R}^{\dinX\times\rpad},\\
d\Bpad &=\Umix^{-\mathsf T}d\Bt\Nout^{\mathsf T}
          \in\mathbb{R}^{\rpad\times\doutX},
\label{eq:method-lora-gradient-recovery}
\end{align}
and adds the controller-computed pad-gradient compensation when $\Tpad$ was
used. It extracts the first $\rtrue$ columns of $d\Apad$ and the first
$\rtrue$ rows of $d\Bpad$ as $d\Amat$ and $d\Bmat$.

The optimizer interface accepts the controller-owned factors, their recovered
real-rank gradients, and controller-owned optimizer state. SGD or momentum may
instead consume the corresponding masked factors and masked buffers when all
participating masks use the supported orthogonal family. Coordinatewise
second-moment optimizers use the trusted interface: $\ctrl$ updates the
plaintext factors and optimizer state, constructs fresh padded factors and
rank masks, and sends their new transformed images to $\acc$. A training step
therefore transitions from private factors, to transformed forward state, to
transformed backward state, back to private gradients and optimizer state, and
finally to refreshed transformed factors. At training termination, $\ctrl$
retains the real-rank adapter and erases the pass-local padded and transformed
states.

\subsection{Unified Execution Lifecycle}
\label{subsec:unified-lifecycle}

Algorithm~\ref{alg:unified-lifecycle} is the implementation reference joining
initialization, prefill, decode, and optional LoRA training. A deployment may
invoke only inference or only adaptation, but it keeps the same ownership and
boundary rules.

\begin{methodalgorithm}{Unified inference and LoRA-adaptation lifecycle}
\label{alg:unified-lifecycle}
\textbf{Inputs:} private $x$ and $\params$; public $\arch$; optional private
training data and $(\Amat,\Bmat)$.\par
\textbf{Outputs:} generated tokens and/or updated private $(\Amat,\Bmat)$.\par
\textbf{Trusted state:} $\params,\maskset,\Tpad,\Recover$; plaintext hidden
state; session cache masks; recovered logits; optional gradients and optimizer
state.\par
\textbf{Accelerator state:} $\Wt$ and optional $\At,\Bt$; current $\Xt,\Yt$;
per-layer $\Kt,\Vt$.\par
\textbf{Communication:} \textsc{Send} denotes $\ctrl\rightarrow\acc$ and
\textsc{Return} denotes $\acc\rightarrow\ctrl$.\par
\begin{enumerate}
\setlength{\itemsep}{0pt}\setlength{\parskip}{0pt}\setlength{\parsep}{0pt}
\item \textbf{Initialize:} $\ctrl$ installs $\params$ and optional adapter;
samples session $\NK,\NV$; validates the call plan; and prepares the call-local
mask/pad schedule. It sends only public session metadata to $\acc$.
\item \textbf{Prefill:} $\ctrl$ embeds $x$. For $\ell=1,\ldots,L$, execute
trusted normalization, masked $Q/K/V$ projections, trusted RoPE and head
masking, accelerator attention, masked output projection, trusted residual
update, masked MLP, and trusted residual update.
\item $\acc$ stores the resulting $\Kt,\Vt$ for every layer. It executes the
masked LM head and \textsc{Return}s masked logits. $\ctrl$ recovers logits and
selects one token.
\item \textbf{Decode:} while the stopping rule is false, $\ctrl$ embeds only
the selected token; repeat Step 2 with one new row, appending that row to each
masked cache; recover logits and select the next token.
\item \textbf{Adapt, if requested:} for each training step and adapted linear,
$\ctrl$ prepares $\Xt,\Wt,\At,\Bt,\Ct$ and \textsc{Send}s them; $\acc$
\textsc{Return}s $\Yt$. $\ctrl$ evaluates the trusted loss and constructs
$\Gt$.
\item \textsc{Send} $\Xt,\At,\Bt,\Gt$; $\acc$ computes and
\textsc{Return}s $d\At,d\Bt$. $\ctrl$ recovers $d\Amat,d\Bmat$, updates the
private factors and optimizer state, and refreshes the accelerator images.
\item \textbf{Terminate:} erase plaintext pass state, call-local masks and pads,
session cache masks, and the masked KV cache; retain only the requested output
tokens and/or updated private adapter.
\end{enumerate}
\end{methodalgorithm}
